\subsubsection{Transporte electrónico bidireccional y registro aritmético}
\label{ii:transporte-electronico}

La representación central utiliza una celda compartida por las dos
evaluaciones APP. Sus coordenadas pertenecen a \(D_9=\{1,\ldots,9\}\),
y la inversión \((r,c)\mapsto(10-r,10-c)\) tiene el punto fijo único
\((5,5)\). El estado visible se escribe \((r,c,h,v,n)\), con
\(h,v\in\{-1,1\}\) e índice cronológico entero \(n\geq0\).
El selector trítico repite \((+1,-1,0)\). En la hoja aditiva se lee
\(r+c\) y se traslada la columna en la dirección \(h\); en la
multiplicativa se lee \(rc\) y se traslada la fila en la dirección
\(v\). El umbral conserva la posición. Ambas orientaciones se
invierten al completar cada intervalo de veintisiete actualizaciones.

Para expresar la lectura y el acarreo con el mismo convenio positivo,
definimos
\[
 [z]^+_9=1+((z-1)\bmod9),\qquad
 q_9(z)=\frac{z-[z]^+_9}{9}.
\]
La evaluación precede al desplazamiento. Si \(a_n\) es el entero
leído, el evento conserva \((a_n,[a_n]^+_9,q_9(a_n))\), la hoja,
la posición anterior, la orientación y el cruce de frontera. El
símbolo trítico es \([a_n]^+_9\bmod3\). El levantamiento espacial
retiene, por ejemplo, el entero
\(q_9(c+h)\) cuando \(c\) se sustituye por \([c+h]^+_9\).

La representación de celda compartida y la representación de dos
cursores de \eqref{ii:tpk-fobs-posiciones} tienen estados visibles
diferentes. Sus registros quedan relacionados con el estado TPK por
sus mapas de representación respectivos. Para la trayectoria central
se utiliza aquí, de forma íntegra, la actualización sobre la celda
compartida que produce la palabra electrónica.

\begin{proposition}[Actualización invertible con cocientes conservados]
\label{ii:electron-inversa-cocientes}
Añadamos al estado visible dos registros enteros \(Q_+,Q_\times\)
y dos números de enrollamiento \(W_r,W_c\). Cada lectura activa
incrementa el registro de su hoja en \(q_9(a_n)\); cada desplazamiento
incrementa el enrollamiento correspondiente en su cruce de frontera.
Esta actualización posee una inversa explícita sobre su imagen.
\end{proposition}
\begin{proof}
El índice posterior determina el modo anterior y si acaba de ocurrir
una inversión de orientaciones. Se deshace primero esa inversión.
Si el modo era aditivo, se recupera la columna previa como
\([c-h]^+_9\); si era multiplicativo, se recupera la fila como
\([r-v]^+_9\). El modo neutro conserva ambas coordenadas. La celda
recuperada determina nuevamente el entero leído, su residuo, su
cociente y el cruce de frontera. Restando esos incrementos de
\(Q_+,Q_\times,W_r,W_c\) e invirtiendo el índice cronológico se
recupera exactamente el estado precedente. La composición de estas
inversas recupera una historia finita completa. La concatenación de
eventos conserva además el orden de las lecturas.
\end{proof}

Los registros \(Q_+,Q_\times\) son una representación mediante sumas explícitas
de la historia de cocientes. La construcción se integra en la fibra
de memoria junto con sus demás coordenadas. Su ley de composición,
para historias concatenables \(u,v\), es
\[
 Q_s(uv)=Q_s(u)+Q_s(v),\qquad s\in\{+,\times\},
\]
donde la segunda historia comienza en el estado final de la primera.
La inversión algebraica de una actualización resta sus incrementos;
recorrer posteriormente un trayecto espacial de vuelta registra
nuevas evaluaciones. Esa distinción conserva simultáneamente la
reversibilidad del transporte y la duración de la historia.

\begin{proposition}[Retorno central y acumulación exacta]
\label{ii:electron-memoria-retorno}
Desde \((r,c,h,v,n)=(5,5,1,1,0)\) y registros iniciales nulos, la
proyección espacial y sus orientaciones retornan después de cincuenta y cuatro
actualizaciones. En ese intervalo,
\[
 (\Delta W_r,\Delta W_c)=(0,0),\qquad
 (\Delta Q_+,\Delta Q_\times)=(10,46).
\]
Después de \(54N\) actualizaciones, para todo entero \(N\geq0\),
se tiene \((Q_+,Q_\times)=(10N,46N)\), mientras las coordenadas
espaciales y sus orientaciones retornan. El índice entero toma el valor
\(54N\); el calendario dodecafásico módulo \(108\) retorna para
\(N\) par.
\end{proposition}
\begin{proof}
Cada par de modos activos desplaza una vez cada coordenada. En
veintisiete actualizaciones ambas recorren una vuelta de nueve
posiciones y retornan a la diagonal central; después se invierten
\(h,v\). El bloque siguiente recorre la vuelta opuesta. Los cruces
espaciales se cancelan entre ambos bloques.

En cada bloque las lecturas aditivas recorren \(2k\),
\(1\leq k\leq9\), cuyo cociente positivo suma cinco. Las lecturas
multiplicativas del bloque positivo recorren
\(k[k+1]^+_9\). Su suma entera y su suma de residuos son
\[
 \sum_{k=1}^9 k[k+1]^+_9=249,\qquad
 \sum_{k=1}^9[k[k+1]^+_9]^+_9=42.
\]
La segunda identidad resulta de la lista
\((2,6,3,2,3,6,2,9,9)\). Por división euclídea,
la suma de cocientes es \((249-42)/9=23\). La sustitución
\(k\mapsto[k-1]^+_9\) muestra que el bloque negativo contiene
el mismo multiconjunto de productos. Ambos bloques producen, por
tanto, diez y cuarenta y seis unidades de los registros respectivos.
La proyección espacial y sus orientaciones vuelven a sus valores
iniciales. Las reglas de activación y de inversión se repiten al
desplazar el índice en \(54\), y los incrementos no dependen de los
registros acumulados. La ley aditiva
anterior prueba por inducción la fórmula para todo \(N\).
\end{proof}

La lectura trítica y la orientación por ventanas proporcionan, sobre
la misma trayectoria,
\[
 \gamma_e=\bigl((100000220)^3(120200000)^3\bigr)^2,\qquad
 \tau_e=(+,+,+,-,-,-,+,+,+,-,-,-).
\]
Las potencias representan concatenaciones. Los lectores de
discrepancias y de ventanas desarrollados en
\eqref{eq:el-kd} y \eqref{eq:el-komega} calculan ambos
\((80,54,6)\). La evaluación se obtiene de la trayectoria antes
de incorporar \(\alpha\), la sección de acción o una masa.
La memoria de cocientes conserva información adicional a ese triple:
en ciento ocho actualizaciones registra \((20,92)\), con
enrollamiento espacial neto nulo. Cada representación mantiene así el
tipo del observable que la produce.
